\documentclass[12pt]{article}
\usepackage[utf8]{inputenc}
\usepackage[T1]{fontenc}
\usepackage{hyperref}
\title{Kuku-Mangk Language Resources for K-12 Curriculum Development}
\author{Compiled by Claude Sonnet 5.5 (claude-sonnet-5-5) for Andr\'as Kornai}
\date{}
\begin{document}
\maketitle

\section{Introduction}
Kuku-Mangk (ISO 639-3 xmq, Glottocode kuku1279) is an extinct Australian language, apparently from Cape York Peninsula. The searches found no dedicated Kuku-Mangk grammar, dictionary, primer or audio. One Wikipedia page lists the xmq code as unattested. Please do not confuse this language with Kuku Yalanji, Kuuku Ya'u or Kugu Nganhcara, which are different languages with their own materials. A Kuku-Mangk curriculum would have to be built from archival sources that have not yet been located.

\section{Primary Educational Resources}
No primers, textbooks or school programmes for Kuku-Mangk were found.

\section{Linguistic and Reference Materials}
\begin{itemize}
\item Sidney H. Ray (1907), a short comparative vocabulary of the languages of Cape York Peninsula, is catalogued in Glottolog. The search did not confirm that it contains Kuku-Mangk material. \url{https://glottolog.org/resource/reference/id/472081}
\item The Glosbe page labelled ``Kuku-Mangk'' is a mislabelled pairing with Metlat\'onoc Mixtec and looks unreliable. \url{https://glosbe.com/mxv/xmq}
\item For context only, \emph{A Dictionary of Umpithamu, with notes on Middle Paman} documents a neighbouring Cape York language, not Kuku-Mangk. \url{https://shop.slq.qld.gov.au/products/9781925302219}
\end{itemize}

\section{Religious and Community Resources}
The Pama Language Centre covers Cape York languages. Its Kuuku Ya'u page is a different language, and it states that no dictionary is currently available for that variety. No Kuku-Mangk community organisation was found. \url{https://www.pamacentre.org.au/kuuku-yau/}

\section{Digital and Audio Resources}
No Kuku-Mangk audio, video or software was found.

\section{Curriculum Development Context}
Little or nothing is publicly documented for Kuku-Mangk. Any heritage curriculum would first need archival work, for example with AIATSIS or the Pama Language Centre, and consultation with Traditional Owner descendants. The search results do not confirm that either holds Kuku-Mangk material.
\end{document}